\begin{abstract}
多模态情感分析需要同时处理文本、音频和视觉序列的长度差异、局部缺失和情感标签的不确定性。本文建立了一套可复核的多模态时序建模流程，包含特征提取与统一长度处理、缺失模态鲁棒训练、三分类与连续情感强度联合预测，以及基于遮挡敏感性的结果解释。

针对问题一，本文从文本、音频和视觉三种模态提取特征，以文本有效序列为基准，将音频和视觉按时间顺序分段聚合，统一为最多 $L=50$ 个时序位置；不足位置进行零填充，并保存有效长度和缺失掩码，形成从原始样本到统一多模态张量的可复现映射。

针对问题二，本文构建 DRM-Net。各模态特征首先投影到统一空间，再由 Bi-GRU 提取模态内时序信息；跨模态注意力用于利用可用模态进行补偿，动态门控根据输入状态分配模态权重，最后联合输出三分类极性和连续情感强度。训练阶段模拟局部连续缺失，并显式传入缺失掩码。模型选择只使用训练集和验证集，最终测试集仅在方案锁定后评估。三个固定随机种子的最佳验证模型在测试时进行分类概率平均并采用 Argmax，最终得到 Accuracy $66.30\%$、Macro-F1 $62.54\%$、MAE $0.6346$ 和 Pearson $r=0.6630$。结构消融显示，缺失增强和缺失掩码对分类性能具有明显作用。附件 3 为无标签专项测试集，仅报告预测结果，不计算准确率或 F1。

针对问题三，本文采用两层解释流程：首先分别遮挡文本、音频和视觉输入，计算预测变化并得到模态遮挡敏感性；随后在主要模态内逐时间步遮挡，定位显著时间窗口，并映射到视频时间、帧区间和文本片段。该解释依赖已训练模型且受模态交互影响，因此不能解释为严格独立因果效应。附件 4 共输出 20 条无标签预测与解释记录。

\keywords{多模态情感分析\quad 时序统一\quad 缺失模态鲁棒性\quad 动态门控\quad 遮挡敏感性\quad 可解释性}
\end{abstract}
